% LaTeX Template for AI Problem Analysis Output (v2.0)
% AMC12B 2023 Problem 11 Strategic Analysis
% Generated using AMC12_COMC_Problem_Solving_Prompt.md framework

\documentclass[12pt]{article}
\usepackage[utf8]{inputenc}
\usepackage{amsmath, amssymb, amsthm}
\usepackage{geometry}
\usepackage{xcolor}
\usepackage[most]{tcolorbox}
\usepackage{enumitem}
\usepackage{fancyhdr}

% Page setup
\geometry{margin=1in}
\setlength{\headheight}{14.5pt}
\pagestyle{fancy}
\fancyhf{}
\rhead{AMC12B 2023 Problem 11 Analysis}
\lhead{\today}

% Color definitions
\definecolor{problemconfig}{RGB}{230, 242, 255}    % Light blue
\definecolor{strategic}{RGB}{255, 230, 242}       % Light pink
\definecolor{tactical}{RGB}{242, 255, 230}        % Light green
\definecolor{techniques}{RGB}{255, 242, 230}      % Light yellow
\definecolor{verification}{RGB}{255, 230, 230}    % Light red
\definecolor{solution}{RGB}{230, 255, 255}        % Light cyan
\definecolor{competition}{RGB}{242, 242, 255}     % Light purple
\definecolor{proofwriting}{RGB}{240, 230, 255}    % Light lavender
\definecolor{gold}{RGB}{218, 165, 32}

% Box styles
\newtcolorbox{problemconfigbox}{
    colback=problemconfig,
    colframe=blue,
    title=Problem Configuration Analysis,
    fonttitle=\bfseries,
    arc=3pt,
    boxrule=1pt,
    breakable
}

\newtcolorbox{strategicbox}{
    colback=strategic,
    colframe=purple,
    title=Strategic Framework Application,
    fonttitle=\bfseries,
    arc=3pt,
    boxrule=1pt,
    breakable
}

\newtcolorbox{tacticalbox}{
    colback=tactical,
    colframe=green,
    title=Tactical Methodology Selection,
    fonttitle=\bfseries,
    arc=3pt,
    boxrule=1pt,
    breakable
}

\newtcolorbox{techniquesbox}{
    colback=techniques,
    colframe=orange,
    title=Conversion Techniques Application,
    fonttitle=\bfseries,
    arc=3pt,
    boxrule=1pt,
    breakable
}

\newtcolorbox{verificationbox}{
    colback=verification,
    colframe=red,
    title=Verification Strategy Design,
    fonttitle=\bfseries,
    arc=3pt,
    boxrule=1pt,
    breakable
}

\newtcolorbox{solutionbox}{
    colback=solution,
    colframe=cyan,
    title=Solution Roadmap,
    fonttitle=\bfseries,
    arc=3pt,
    boxrule=1pt,
    breakable
}

\newtcolorbox{competitionbox}{
    colback=competition,
    colframe=purple,
    title=Competition Strategy Integration,
    fonttitle=\bfseries,
    arc=3pt,
    boxrule=1pt,
    breakable
}

\newtcolorbox{keyinsightbox}{
    colback=yellow!20,
    colframe=gold,
    title=Key Insight,
    fonttitle=\bfseries,
    arc=3pt,
    boxrule=2pt,
    leftrule=5pt,
    breakable
}

\newtcolorbox{warningbox}{
    colback=red!10,
    colframe=red!50,
    title=Warning: Common Pitfall,
    fonttitle=\bfseries,
    arc=3pt,
    boxrule=2pt,
    leftrule=5pt,
    breakable
}

\newtcolorbox{tipbox}{
    colback=green!10,
    colframe=green!50,
    title=Pro Tip,
    fonttitle=\bfseries,
    arc=3pt,
    boxrule=2pt,
    leftrule=5pt,
    breakable
}

\begin{document}

\title{AMC12B 2023 Problem 11\\Strategic Problem Analysis}
\author{Competition Math Framework v2.1}
\date{\today}
\maketitle

\section*{Problem Statement}

What is the maximum area of an isosceles trapezoid that has legs of length $1$ and one base twice as long as the other?

$$\textbf{(A) }\frac{5}{4} \qquad \textbf{(B) } \frac{8}{7} \qquad \textbf{(C)} \frac{5\sqrt{2}}{4} \qquad \textbf{(D) } \frac{3}{2} \qquad \textbf{(E) } \frac{3\sqrt{3}}{4}$$

\begin{problemconfigbox}
\subsection*{A. Problem Classification}
\begin{itemize}[leftmargin=*]
    \item \textbf{Problem Type}: To Find (maximum value)
    \item \textbf{Competition Context}: AMC12B 2023, Problem 11
    \item \textbf{Difficulty Band}: Mid-band (P11) - Intermediate level, 3-5 minutes expected
    \item \textbf{Mathematical Domain}: Geometry (Optimization, Trapezoid Geometry) with Algebraic Techniques
    \item \textbf{Estimated Time}: 3-5 minutes with verification
\end{itemize}

\subsection*{B. Problem Structure Identification}
\begin{itemize}[leftmargin=*]
    \item \textbf{Given Information}: 
    \begin{itemize}
        \item Isosceles trapezoid (two legs of equal length)
        \item Legs have length $1$
        \item One base is twice as long as the other base
        \item Need to find maximum possible area
    \end{itemize}
    \item \textbf{Constraints}: 
    \begin{itemize}
        \item Legs must be length $1$ (fixed)
        \item Base ratio must be $2:1$ (fixed relationship)
        \item Must form a valid trapezoid (geometric feasibility)
        \item Isosceles symmetry condition
    \end{itemize}
    \item \textbf{Objective}: Find the maximum area achievable under these constraints
    \item \textbf{Hidden Assumptions}: 
    \begin{itemize}
        \item The trapezoid must be convex
        \item Bases are parallel (definition of trapezoid)
        \item Variable parameter: the length of the shorter base (which determines everything else)
        \item Geometric feasibility: shorter base $x$ must satisfy $0 < x < 2$
    \end{itemize}
    \item \textbf{Answer Choice Leverage}: 
    \begin{itemize}
        \item Mix of simple fractions and expressions with radicals
        \item Option (D) $\frac{3}{2} = 1.5$ is clean and relatively large
        \item Options with $\sqrt{2}$ and $\sqrt{3}$ suggest optimization might yield irrational critical points
        \item Can estimate: reasonable trapezoid area should be between $0.5$ and $2$
        \item All options are positive and less than $2$
    \end{itemize}
\end{itemize}

\subsection*{C. Strategic Orientation}
\begin{itemize}[leftmargin=*]
    \item \textbf{Hypothesis}: Given fixed leg lengths and base ratio, find optimal base length
    \item \textbf{Conclusion}: Determine the maximum area value
    \item \textbf{Notation}: 
    \begin{itemize}
        \item Let shorter base = $x$, then longer base = $2x$
        \item Height of trapezoid = $h$ (to be expressed in terms of $x$)
        \item Area formula: $A = \frac{1}{2}(b_1 + b_2) \cdot h = \frac{1}{2}(x + 2x) \cdot h = \frac{3xh}{2}$
        \item Pythagorean relation connects $x$ and $h$ via leg length $1$
    \end{itemize}
    \item \textbf{Intuitive Approach}: This is an optimization problem. Express area in terms of one variable, then maximize using algebra (AM-GM) or calculus
    \item \textbf{Cross-Domain Potential}: 
    \begin{itemize}
        \item Pure geometry: Similar triangles, height formula via Pythagorean theorem
        \item Algebraic optimization: AM-GM inequality application
        \item Calculus: Derivative to find critical points (if familiar with calculus)
        \item Trigonometry: Parametric approach using angles
    \end{itemize}
\end{itemize}
\end{problemconfigbox}

\begin{strategicbox}
\subsection*{A. Psychological Strategies}
\begin{itemize}[leftmargin=*]
    \item \textbf{Mental Toughness}: This is an optimization problem requiring algebraic manipulation. Don't be discouraged if the area expression looks messy initially - it simplifies nicely for AM-GM application.
    \item \textbf{Creativity Cultivation}: Multiple valid approaches exist:
    \begin{itemize}
        \item AM-GM inequality (elegant, competition-standard)
        \item Calculus derivative (if you know calculus)
        \item Similar triangles (geometric insight)
        \item Trigonometric substitution (clean but less common)
    \end{itemize}
    \item \textbf{Iterative Refinement}: If your area formula doesn't simplify well, check the geometry setup. The key is properly relating height to base via the Pythagorean theorem.
\end{itemize}

\subsection*{B. Investigation Strategies}
\begin{itemize}[leftmargin=*]
    \item \textbf{Startup Orientation}: 
    \begin{enumerate}
        \item Draw an isosceles trapezoid with labeled dimensions
        \item Identify the variable (shorter base length $x$)
        \item Set up coordinate system or drop perpendiculars to find height
    \end{enumerate}
    \item \textbf{Get Your Hands Dirty}: Draw the trapezoid clearly:
    \begin{itemize}
        \item Place longer base $2x$ horizontally
        \item Shorter base $x$ parallel above it
        \item By symmetry, each leg projects horizontally by $\frac{x}{2}$
        \item Height $h$ and horizontal projection satisfy: $h^2 + \left(\frac{x}{2}\right)^2 = 1^2$
    \end{itemize}
    \item \textbf{Penultimate Step}: After setting up area expression $A(x)$, the final step is applying optimization technique (AM-GM or derivative) to find maximum.
    \item \textbf{Make It Easier}: Instead of general trapezoid, focus on the specific constraint (base ratio $2:1$). This reduces complexity significantly.
    \item \textbf{Conjecture via Small Cases}: 
    \begin{itemize}
        \item If $x$ is very small: trapezoid is nearly a line segment, area $\approx 0$
        \item If $x$ approaches $2$: height approaches $0$, area $\approx 0$
        \item Maximum must occur at some intermediate value
    \end{itemize}
    \item \textbf{Visualization}: Sketch helps tremendously:
    \begin{itemize}
        \item Visualize how changing $x$ affects both base sum and height
        \item As $x$ increases, base sum increases but height decreases
        \item Optimal balance occurs at specific $x$ value
    \end{itemize}
\end{itemize}

\subsection*{C. Argument Construction}
\begin{itemize}[leftmargin=*]
    \item \textbf{Proof Method}: Direct computation with optimization
    \item \textbf{Key Lemmas}: 
    \begin{itemize}
        \item Height formula: $h = \sqrt{1 - \frac{x^2}{4}}$ (from Pythagorean theorem)
        \item Area formula: $A = \frac{3x}{2} \sqrt{1 - \frac{x^2}{4}}$
        \item AM-GM inequality: For positive $a, b$: $ab \leq \frac{(a+b)^2}{4}$ with equality when $a = b$
    \end{itemize}
    \item \textbf{Logical Structure}: 
    \begin{enumerate}
        \item Express height in terms of base length
        \item Substitute into area formula
        \item Apply optimization technique (AM-GM preferred)
        \item Solve for optimal $x$ and compute maximum area
    \end{enumerate}
\end{itemize}
\end{strategicbox}

\begin{tacticalbox}
\subsection*{A. Core Mathematical Tactics}
\begin{itemize}[leftmargin=*]
    \item \textbf{Primary Tactics}: 
    \begin{itemize}
        \item \textbf{Optimization via AM-GM}: Convert area expression to form suitable for AM-GM application
        \item \textbf{Pythagorean Theorem}: Relate height to base length via right triangle formed by dropping perpendicular
        \item \textbf{Symmetry Exploitation}: Isosceles trapezoid has vertical line of symmetry
    \end{itemize}
    \item \textbf{Domain-Specific Tactics}: 
    \begin{itemize}
        \item \textbf{Trapezoid Properties}: Area = $\frac{1}{2}$(sum of bases)(height)
        \item \textbf{Similar Triangles}: Can extend legs to form larger triangle containing trapezoid
        \item \textbf{Coordinate Geometry}: Can place trapezoid on coordinate system for algebraic approach
    \end{itemize}
    \item \textbf{Crossover Tactics}: 
    \begin{itemize}
        \item \textbf{Geometric-Algebraic}: Express geometric constraint algebraically, then optimize
        \item \textbf{Trigonometric Substitution}: Use parametric form $x = 2\cos\theta$ for elegant solution
    \end{itemize}
\end{itemize}
\end{tacticalbox}

\begin{techniquesbox}
\subsection*{A. High-Probability Techniques}
\begin{itemize}[leftmargin=*]
    \item \textbf{Primary Technique}: \textbf{AM-GM Inequality} (H7 in conversion techniques)
    \item \textbf{Recognition Cues}: 
    \begin{itemize}
        \item Optimization problem (finding maximum)
        \item Expression involves product of terms
        \item AMC12 mid-band level
        \item Area formula will have form that fits AM-GM pattern
    \end{itemize}
    \item \textbf{Application}: 
    \begin{enumerate}
        \item Derive area: $A = \frac{3x}{2}\sqrt{1 - \frac{x^2}{4}} = \frac{3}{4}\sqrt{x^2(4-x^2)}$
        \item Recognize AM-GM form: need to maximize $\sqrt{x^2(4-x^2)}$
        \item Apply AM-GM: $x^2(4-x^2) \leq \left(\frac{x^2 + (4-x^2)}{2}\right)^2 = 4$
        \item Equality when $x^2 = 4-x^2$, so $x = \sqrt{2}$
        \item Maximum area: $A = \frac{3}{4} \cdot 2 = \frac{3}{2}$
    \end{enumerate}
\end{itemize}

\subsection*{B. Alternative Techniques}
\begin{itemize}[leftmargin=*]
    \item \textbf{Backup Technique 1}: \textbf{Calculus - Derivative Method}
    \begin{itemize}
        \item Express area: $A(x) = \frac{3x}{2}\sqrt{1 - \frac{x^2}{4}}$
        \item Take derivative: $\frac{dA}{dx} = \frac{3}{2}\sqrt{1-\frac{x^2}{4}} - \frac{3x^2}{4\sqrt{1-\frac{x^2}{4}}}$
        \item Set equal to zero and solve for critical point
        \item Find $x = \sqrt{2}$, evaluate $A(\sqrt{2}) = \frac{3}{2}$
    \end{itemize}
    
    \item \textbf{Backup Technique 2}: \textbf{Similar Triangles (Geometric)}
    \begin{itemize}
        \item Extend trapezoid legs to meet at point $E$ above
        \item Form similar triangles with ratio $1:2$
        \item Larger triangle has legs of length $2$
        \item Maximize isosceles triangle area (achieved when angle = $90^\circ$)
        \item Use ratio to find trapezoid area
    \end{itemize}
    
    \item \textbf{Backup Technique 3}: \textbf{Trigonometric Parametrization}
    \begin{itemize}
        \item Use constraint $x^2/4 + h^2 = 1$ to set $x/2 = \cos\theta$, $h = \sin\theta$
        \item Area becomes $A = 3\cos\theta \sin\theta = \frac{3}{2}\sin(2\theta)$
        \item Maximum when $\sin(2\theta) = 1$, giving $A_{max} = \frac{3}{2}$
    \end{itemize}
    
    \item \textbf{Technique Integration}: 
    \begin{itemize}
        \item AM-GM is fastest and most elegant for competition
        \item Calculus provides verification for those who know it
        \item Similar triangles offers pure geometric insight
        \item Trigonometry is cleanest but requires familiarity with parametrization
    \end{itemize}
\end{itemize}
\end{techniquesbox}

\begin{verificationbox}
\subsection*{A. Verification Level Selection}
\begin{itemize}[leftmargin=*]
    \item \textbf{Chosen Level}: Level 4 (Standard Verification)
    \item \textbf{Justification}: 
    \begin{itemize}
        \item Mid-band problem (P11) with multi-step algebra
        \item Optimization problem where errors in setup are common
        \item Expected confidence: 75-85\% after first solve
        \item Should verify: height formula, area expression, AM-GM application
    \end{itemize}
    \item \textbf{Time Budget}: 45-60 seconds for verification
\end{itemize}

\subsection*{B. Verification Checklist}
\begin{itemize}[leftmargin=*]
    \item \textbf{Domain Restrictions}: 
    \begin{itemize}
        \item[$\checkmark$] Verify $x = \sqrt{2}$ gives valid trapezoid: $0 < \sqrt{2} < 2$ $\checkmark$
        \item[$\checkmark$] Height at $x = \sqrt{2}$: $h = \sqrt{1 - \frac{2}{4}} = \sqrt{\frac{1}{2}} = \frac{1}{\sqrt{2}}$ $\checkmark$
        \item[$\checkmark$] Check geometric feasibility: positive height and bases $\checkmark$
    \end{itemize}
    \item \textbf{Computational Accuracy}: 
    \begin{itemize}
        \item[$\checkmark$] Height formula: $h = \sqrt{1 - \frac{x^2}{4}}$ from $h^2 + \left(\frac{x}{2}\right)^2 = 1$ $\checkmark$
        \item[$\checkmark$] Area formula: $A = \frac{3x}{2} \cdot h = \frac{3x}{2}\sqrt{1-\frac{x^2}{4}}$ $\checkmark$
        \item[$\checkmark$] Simplification: $A = \frac{3}{4}\sqrt{x^2(4-x^2)}$ $\checkmark$
        \item[$\checkmark$] AM-GM application: $x^2(4-x^2) \leq 4$ with equality at $x^2 = 2$ $\checkmark$
        \item[$\checkmark$] Final calculation: $A_{max} = \frac{3}{4}\sqrt{4} = \frac{3}{4} \cdot 2 = \frac{3}{2}$ $\checkmark$
    \end{itemize}
    \item \textbf{Alternative Method Verification}:
    \begin{itemize}
        \item Using $x = \sqrt{2}$: bases are $\sqrt{2}$ and $2\sqrt{2}$, sum = $3\sqrt{2}$
        \item Height: $h = \frac{1}{\sqrt{2}}$
        \item Area: $\frac{1}{2} \cdot 3\sqrt{2} \cdot \frac{1}{\sqrt{2}} = \frac{3}{2}$ $\checkmark$
    \end{itemize}
    \item \textbf{Edge Cases}: 
    \begin{itemize}
        \item[$\checkmark$] Boundary behavior: As $x \to 0$ or $x \to 2$, area $\to 0$ $\checkmark$
        \item[$\checkmark$] Answer matches choice (D): $\frac{3}{2}$ $\checkmark$
        \item[$\checkmark$] Reasonableness: $1.5$ is plausible for such trapezoid $\checkmark$
    \end{itemize}
\end{itemize}
\end{verificationbox}

\begin{solutionbox}
\subsection*{A. Step-by-Step Solution}

\subsubsection*{Method 1: AM-GM Inequality (Recommended - 3 minutes)}

\begin{enumerate}[leftmargin=*]
    \item \textbf{Setup - Find Height Formula (45 seconds)}
    
    Draw an isosceles trapezoid with:
    \begin{itemize}
        \item Shorter base (top) = $x$
        \item Longer base (bottom) = $2x$
        \item Both legs = $1$
    \end{itemize}
    
    Drop perpendiculars from the ends of the top base to the bottom base. By symmetry, each leg projects horizontally a distance of $\frac{2x - x}{2} = \frac{x}{2}$.
    
    By the Pythagorean theorem on the right triangle formed:
    \begin{equation}
        h^2 + \left(\frac{x}{2}\right)^2 = 1^2
    \end{equation}
    
    Solving for height:
    \begin{equation}
        h = \sqrt{1 - \frac{x^2}{4}}
    \end{equation}
    
    \item \textbf{Express Area Formula (30 seconds)}
    
    Using the trapezoid area formula:
    \begin{align*}
        A &= \frac{1}{2}(b_1 + b_2) \cdot h \\
        &= \frac{1}{2}(x + 2x) \cdot \sqrt{1 - \frac{x^2}{4}} \\
        &= \frac{3x}{2} \sqrt{1 - \frac{x^2}{4}}
    \end{align*}
    
    Factor out constants:
    \begin{align*}
        A &= \frac{3x}{2} \cdot \frac{1}{2}\sqrt{4 - x^2} \\
        &= \frac{3}{4} \sqrt{x^2(4 - x^2)}
    \end{align*}
    
    \item \textbf{Apply AM-GM Inequality (60 seconds)}
    
    To maximize $A$, we need to maximize $\sqrt{x^2(4-x^2)}$, which is equivalent to maximizing $x^2(4-x^2)$.
    
    By the AM-GM inequality, for positive values $a$ and $b$:
    \begin{equation}
        ab \leq \left(\frac{a+b}{2}\right)^2
    \end{equation}
    
    Let $a = x^2$ and $b = 4-x^2$. Then $a + b = 4$ (constant), so:
    \begin{align*}
        x^2(4-x^2) &\leq \left(\frac{x^2 + (4-x^2)}{2}\right)^2 \\
        &= \left(\frac{4}{2}\right)^2 \\
        &= 4
    \end{align*}
    
    Equality holds when $x^2 = 4-x^2$, which gives:
    \begin{align*}
        2x^2 &= 4 \\
        x^2 &= 2 \\
        x &= \sqrt{2} \quad \text{(taking positive root)}
    \end{align*}
    
    \item \textbf{Calculate Maximum Area (30 seconds)}
    
    \begin{align*}
        A_{max} &= \frac{3}{4}\sqrt{x^2(4-x^2)} \\
        &= \frac{3}{4}\sqrt{4} \\
        &= \frac{3}{4} \cdot 2 \\
        &= \frac{3}{2}
    \end{align*}
    
    \item \textbf{Verification (30 seconds)}
    
    Check with $x = \sqrt{2}$:
    \begin{itemize}
        \item Bases: $\sqrt{2}$ and $2\sqrt{2}$, sum = $3\sqrt{2}$
        \item Height: $h = \sqrt{1-\frac{2}{4}} = \frac{1}{\sqrt{2}}$
        \item Area: $\frac{1}{2} \cdot 3\sqrt{2} \cdot \frac{1}{\sqrt{2}} = \frac{3}{2}$ $\checkmark$
    \end{itemize}
    
    \item \textbf{Final Answer}
    
    The maximum area is $\frac{3}{2}$, which is \textbf{answer choice (D)}.
\end{enumerate}

\subsubsection*{Method 2: Similar Triangles (Geometric - 3-4 minutes)}

\begin{enumerate}[leftmargin=*]
    \item \textbf{Extend Trapezoid to Triangle}
    
    Let trapezoid $ABCD$ have $AB = x$ (shorter base), $CD = 2x$ (longer base), and $AD = BC = 1$ (legs).
    
    Extend legs $AD$ and $BC$ upward until they meet at point $E$.
    
    \item \textbf{Identify Similar Triangles}
    
    Triangles $\triangle ABE$ and $\triangle DCE$ are similar with ratio $1:2$ (since $AB:CD = x:2x = 1:2$).
    
    Since the ratio is $1:2$, and $AD = 1$, we have:
    \begin{itemize}
        \item $ED = 2 \cdot 1 = 2$
        \item $AE = ED - AD = 2 - 1 = 1$
    \end{itemize}
    
    Similarly, $BE = 1$.
    
    \item \textbf{Relate Areas}
    
    Since the linear ratio is $1:2$, the area ratio is $1:4$:
    \begin{equation}
        [DCE] = 4 \cdot [ABE]
    \end{equation}
    
    Therefore:
    \begin{align*}
        [ABCD] &= [DCE] - [ABE] \\
        &= 4[ABE] - [ABE] \\
        &= 3[ABE]
    \end{align*}
    
    Or equivalently: $[ABCD] = \frac{3}{4}[DCE]$
    
    \item \textbf{Maximize Triangle Area}
    
    We need to maximize $[DCE]$, which is an isosceles triangle with legs $ED = EC = 2$.
    
    For an isosceles triangle with legs of length $\ell$ and angle $\theta$ between them:
    \begin{equation}
        \text{Area} = \frac{1}{2}\ell^2 \sin\theta
    \end{equation}
    
    For $\triangle DCE$ with $\ell = 2$:
    \begin{equation}
        [DCE] = \frac{1}{2}(2)(2)\sin\theta = 2\sin\theta
    \end{equation}
    
    This is maximized when $\sin\theta = 1$, i.e., $\theta = 90^\circ$.
    
    Maximum: $[DCE]_{max} = 2(1) = 2$
    
    \item \textbf{Calculate Trapezoid Area}
    
    \begin{align*}
        [ABCD]_{max} &= \frac{3}{4}[DCE]_{max} \\
        &= \frac{3}{4} \cdot 2 \\
        &= \frac{3}{2}
    \end{align*}
\end{enumerate}

\subsubsection*{Method 3: Trigonometric Parametrization (Alternative - 3 minutes)}

\begin{enumerate}[leftmargin=*]
    \item \textbf{Parametrize Using Angle}
    
    From the constraint $h^2 + \left(\frac{x}{2}\right)^2 = 1$, we can write:
    \begin{align*}
        \frac{x}{2} &= \cos\theta \\
        h &= \sin\theta
    \end{align*}
    
    where $0 < \theta < \frac{\pi}{2}$ (first quadrant for positive values).
    
    \item \textbf{Express Area in Terms of $\theta$}
    
    \begin{align*}
        A &= \frac{1}{2}(x + 2x) \cdot h \\
        &= \frac{3x}{2} \cdot h \\
        &= \frac{3(2\cos\theta)}{2} \cdot \sin\theta \\
        &= 3\cos\theta \sin\theta
    \end{align*}
    
    \item \textbf{Use Double-Angle Formula}
    
    \begin{align*}
        A &= 3\cos\theta \sin\theta \\
        &= \frac{3}{2}(2\cos\theta \sin\theta) \\
        &= \frac{3}{2}\sin(2\theta)
    \end{align*}
    
    \item \textbf{Maximize}
    
    Since $0 < \theta < \frac{\pi}{2}$, we have $0 < 2\theta < \pi$.
    
    The maximum value of $\sin(2\theta)$ in this range is $1$, occurring at $2\theta = \frac{\pi}{2}$, i.e., $\theta = \frac{\pi}{4}$.
    
    \item \textbf{Calculate Maximum}
    
    \begin{equation}
        A_{max} = \frac{3}{2} \cdot 1 = \frac{3}{2}
    \end{equation}
\end{enumerate}

\subsection*{B. Solution Comparison}

\begin{itemize}[leftmargin=*]
    \item \textbf{Method 1 (AM-GM) Advantages}: 
    \begin{itemize}
        \item Standard AMC12 technique for optimization
        \item Algebraic and systematic
        \item No calculus or trigonometry required
        \item Clean and efficient
    \end{itemize}
    \item \textbf{Method 2 (Similar Triangles) Advantages}:
    \begin{itemize}
        \item Pure geometric insight
        \item Reduces to maximizing isosceles triangle
        \item Elegant conceptual approach
        \item Good for geometric thinkers
    \end{itemize}
    \item \textbf{Method 3 (Trigonometry) Advantages}:
    \begin{itemize}
        \item Cleanest algebraic form
        \item Natural parametrization
        \item Requires trigonometric facility
        \item Very efficient if comfortable with technique
    \end{itemize}
    \item \textbf{All methods yield}: $A_{max} = \frac{3}{2}$, confirming answer \textbf{(D)}.
\end{itemize}
\end{solutionbox}

\begin{competitionbox}
\subsection*{A. Time Management}
\begin{itemize}[leftmargin=*]
    \item \textbf{Estimated Time}: 3-5 minutes total (including verification)
    \begin{itemize}
        \item Setup and height formula: 45-60 seconds
        \item Method 1 (AM-GM): 2-2.5 minutes
        \item OR Method 2 (Similar triangles): 2.5-3 minutes
        \item OR Method 3 (Trigonometry): 2-2.5 minutes
        \item Verification: 45-60 seconds
    \end{itemize}
    \item \textbf{Time Checkpoints}: 
    \begin{itemize}
        \item At 1 minute: Should have height formula and area expression
        \item At 2.5 minutes: Should be applying optimization technique
        \item At 4 minutes: Should be verifying answer
        \item If stuck after 5 minutes: Make educated guess and move on
    \end{itemize}
    \item \textbf{Priority Level}: \textbf{High} - Accessible P11 with standard optimization techniques. Complete before attempting P17+.
\end{itemize}

\subsection*{B. Risk Assessment}
\begin{itemize}[leftmargin=*]
    \item \textbf{Confidence Level}: 80-85\% (standard optimization with some algebraic complexity)
    \item \textbf{Abandonment Criteria}: 
    \begin{itemize}
        \item If height formula setup is unclear after 2 minutes
        \item If AM-GM pattern isn't recognized after 4 minutes total
        \item If algebra becomes unmanageably messy (check setup)
    \end{itemize}
    \item \textbf{Flag for Review}: Possibly (if AM-GM application felt uncertain)
\end{itemize}

\subsection*{C. Answer Choice Strategy}
\begin{itemize}[leftmargin=*]
    \item \textbf{Elimination Potential}:
    \begin{itemize}
        \item All choices are positive (good - area must be positive)
        \item Can estimate: with legs = 1, reasonable area is 0.5 to 2
        \item (A) $\frac{5}{4} = 1.25$
        \item (B) $\frac{8}{7} \approx 1.14$
        \item (C) $\frac{5\sqrt{2}}{4} \approx 1.77$
        \item (D) $\frac{3}{2} = 1.50$ $\checkmark$
        \item (E) $\frac{3\sqrt{3}}{4} \approx 1.30$
    \end{itemize}
    \item \textbf{Pattern Recognition}:
    \begin{itemize}
        \item Answer involves factor of $\frac{3}{4}$ times something
        \item This matches our area formula structure
        \item The $\frac{3}{2}$ comes from $\frac{3}{4} \cdot 2$
    \end{itemize}
    \item Our answer $\frac{3}{2}$ matches choice (D) exactly.
\end{itemize}
\end{competitionbox}

\begin{keyinsightbox}
\textbf{Key Insight}: This problem combines geometric constraint with algebraic optimization. The crucial steps are:

\begin{enumerate}
    \item \textbf{Setup}: Use Pythagorean theorem to express height in terms of base: $h = \sqrt{1 - \frac{x^2}{4}}$
    
    \item \textbf{Formula}: Area becomes $A = \frac{3}{4}\sqrt{x^2(4-x^2)}$
    
    \item \textbf{Optimization}: Apply AM-GM to the product $x^2(4-x^2)$ where the sum $x^2 + (4-x^2) = 4$ is constant. This immediately gives maximum of $4$ when $x^2 = 2$.
    
    \item \textbf{Result}: $A_{max} = \frac{3}{4} \cdot 2 = \frac{3}{2}$
\end{enumerate}

The AM-GM pattern: When optimizing $a \cdot b$ where $a + b = k$ (constant), the maximum is $\left(\frac{k}{2}\right)^2$ achieved when $a = b = \frac{k}{2}$.
\end{keyinsightbox}

\begin{warningbox}
\textbf{Warning: Common Pitfalls}

\begin{enumerate}
    \item \textbf{Height Formula Error}: Students often forget the factor of $\frac{1}{2}$ in $\left(\frac{x}{2}\right)^2$. The horizontal projection is half the base difference, not the full difference.
    
    \item \textbf{Area Formula Mistake}: Don't forget the $\frac{1}{2}$ in trapezoid area formula, or that both bases contribute (sum = $3x$, not $2x$).
    
    \item \textbf{AM-GM Misapplication}: AM-GM requires $a + b$ to be constant. Here, $x^2 + (4-x^2) = 4$ is indeed constant, which makes AM-GM applicable.
    
    \item \textbf{Domain Issues}: The value $x = \sqrt{2}$ must satisfy $0 < x < 2$ for a valid trapezoid. Check: $\sqrt{2} \approx 1.41 < 2$ $\checkmark$
    
    \item \textbf{Arithmetic Error}: When computing $\frac{3}{4} \cdot 2$, don't mistakenly get $\frac{3}{8}$ or $\frac{6}{4}$ without simplifying.
    
    \item \textbf{Similar Triangles Confusion}: The area ratio is the square of the linear ratio: $(1:2)$ for lengths gives $(1:4)$ for areas.
\end{enumerate}
\end{warningbox}

\begin{tipbox}
\textbf{Pro Tips for Competition Success}

\begin{enumerate}
    \item \textbf{Pattern Recognition}: "Maximize area with constraints" + "product form" → Think AM-GM immediately. This is a P9-P13 level standard pattern.
    
    \item \textbf{Algebraic Setup}: Always express area as a function of a single variable. Here, choosing $x$ (shorter base) as the variable makes everything clean.
    
    \item \textbf{AM-GM Recognition}: When you see $A = c\sqrt{f(x)(k-f(x))}$ where $k$ is constant, AM-GM applies beautifully.
    
    \item \textbf{Quick Check}: After finding $x = \sqrt{2}$, verify it's in valid domain $(0, 2)$ before proceeding.
    
    \item \textbf{Answer Estimation}: Before detailed work, estimate: "Area should be around $1$ to $1.5$ based on dimensions." This helps catch gross errors.
    
    \item \textbf{Multiple Methods}: If you know calculus or trigonometry, these provide quick verification paths. But AM-GM is competition-standard and should be your first choice.
    
    \item \textbf{Geometric Insight}: The similar triangles method is elegant and worth knowing. It shows that optimization reduces to maximizing an isosceles triangle with fixed leg length (achieved at right angle).
\end{enumerate}
\end{tipbox}

\section*{Final Answer}

Using AM-GM inequality to optimize the area expression, we find that the maximum area is:

\begin{equation*}
\boxed{\textbf{(D) } \frac{3}{2}}
\end{equation*}

\section*{Confidence Assessment}
\begin{itemize}[leftmargin=*]
    \item \textbf{Confidence Level}: 90\% - Standard optimization problem with well-established AM-GM technique. Verified with direct calculation.
    \item \textbf{Verification Applied}: 
    \begin{itemize}
        \item Level 4 (Standard Verification)
        \item Primary method: AM-GM inequality application
        \item Cross-check: Direct substitution of $x = \sqrt{2}$ confirms answer
        \item Reasonableness: Value of $1.5$ is plausible given constraints
        \item Domain check: Optimal $x = \sqrt{2}$ satisfies $0 < x < 2$
    \end{itemize}
    \item \textbf{Flag for Review}: No - High confidence with algebraic verification
    \item \textbf{Time Spent}: Approximately 3 minutes for solution + 1 minute for verification = 4 minutes total (within target range for P11)
\end{itemize}

\section*{Pattern Library Entry}

\textbf{Problem Signature}: Geometric optimization with constraint

\textbf{Recognition Cues}: 
\begin{itemize}
    \item Maximize/minimize area or length
    \item Fixed constraints (leg lengths, base ratio)
    \item Single degree of freedom
    \item Product form in area expression
\end{itemize}

\textbf{Primary Technique}: AM-GM inequality for optimization

\textbf{Key Formula}: For product $ab$ with $a+b = k$ constant, maximum is $\frac{k^2}{4}$ when $a = b = \frac{k}{2}$

\textbf{Alternative Approaches}: 
\begin{itemize}
    \item Calculus (derivative)
    \item Similar triangles (geometric reduction)
    \item Trigonometric parametrization
\end{itemize}

\textbf{Common Traps}: Height formula errors, AM-GM misapplication, arithmetic mistakes

\textbf{Time Estimate}: 3-5 minutes

\textbf{Similar Problems}: AMC12 optimization problems (inscribed/circumscribed figures, constrained maxima), trapezoid geometry, AM-GM applications

\end{document}

